\documentclass[UTF8]{ctexart}
\usepackage{graphicx}
\usepackage{booktabs}
\usepackage{geometry}
\usepackage{array}
\usepackage{amsmath}
\geometry{a4paper, margin=2.5cm}
\title{NIPT 检测时点选择与胎儿异常判定模型}
\author{}
\date{}
\begin{document}
\maketitle
\begin{abstract}
本文围绕 NIPT 检测中的时点选择和胎儿异常判定问题，建立男胎 Y 染色体浓度模型、BMI 分组检测时点优化模型以及女胎异常风险判别模型。首先以孕周、BMI、年龄和 GC 含量解释 Y 染色体浓度，模型拟合得到 $R^2=0.963$，平均绝对误差为 $0.317$ 个百分点，可用于反推浓度达到 $4\%$ 阈值的孕周。其次将检测过早造成的未达标风险与检测过晚风险合并为分组风险函数，得到 BMI 分组推荐时点：<28 组 17.5 周，28-32 组 18.0 周，32-36 组 19.0 周，>=36 组 19.5 周，并通过阈值扰动检验推荐结果的稳定性。最后构造基于染色体 Z 值、X 染色体偏离和 GC 含量的异常风险评分，得到 AUC 为 $0.859$，灵敏度为 $0.586$。结果说明，可解释统计模型、分组优化和判别曲线能够形成一条可复核的 NIPT 决策链；受控数据只用于方法验证，不替代真实临床结论。

\textbf{关键词：}NIPT；检测时点；BMI 分组；风险优化；异常判定
\end{abstract}

\section{问题重述}
无创产前检测通过孕妇外周血中的胎儿游离 DNA 信息评估胎儿染色体风险。对男胎而言，Y 染色体浓度是否达到阈值会影响检测可靠性；对女胎而言，需要综合常染色体和性染色体相关指标识别异常风险。检测时间过早可能造成胎儿浓度不足，检测时间过晚又可能增加临床处置风险，因此时点选择本质上是一个在可靠性和风险之间折中的决策问题。

本文将问题拆分为三层：第一，建立男胎 Y 染色体浓度与孕周、BMI 等因素之间的关系；第二，按 BMI 分组选择推荐检测孕周，并检验阈值变化对推荐结果的影响；第三，对女胎样本建立异常风险评分并评价判别效果。

\begin{figure}[htbp]
\centering
\includegraphics[width=0.86\textwidth]{../figures/fig_problem_overview.png}
\caption{NIPT 检测时点与异常判定的建模对象、输入指标和验证线索}
\end{figure}

\section{问题分析}
Y 染色体浓度一般随孕周增加而上升，但 BMI 较高时胎儿 DNA 相对浓度可能下降，因此单独用孕周判断会忽略个体差异。BMI 分组后，推荐检测时点既要使组内大多数样本达到浓度阈值，又不能无节制推迟。女胎异常判定则不能依赖 Y 染色体浓度，需要从 13、18、21 号染色体 Z 值、X 染色体偏离和 GC 含量等指标构造风险评分。本文采用线性浓度模型、分组网格搜索和 Logistic 判别模型，是因为三者分别对应解释、决策和分类，且每一步都有明确的验证指标。

\section{数据预处理}
对男胎数据，保留孕周、BMI、年龄、GC 含量和 Y 染色体浓度，并以 $4\%$ 作为达标阈值。对女胎数据，构造最大常染色体 Z 值、X 染色体偏离绝对值、GC 偏离量、年龄和 BMI 等特征。所有连续变量在建模前检查范围和缺失情况；分类模型中特征先标准化，以免量纲差异影响权重估计。本文使用受控数据进行流程验证，因此评价指标用于检查模型链路和监督返工，而不是声称真实医学性能。

\section{模型假设}
\begin{enumerate}
\item 在分析窗口内，Y 染色体浓度随孕周的主效应可用线性项近似。
\item BMI、年龄和 GC 含量对浓度的影响在受控样本中可用加性主效应表示。
\item 同一 BMI 组内采用统一推荐检测时点，组内个体差异进入未达标率和敏感性分析。
\item 女胎异常风险可由染色体 Z 值、X 染色体偏离和 GC 偏离的组合评分近似刻画。
\item 判别阈值服务于筛查任务，因此优先保证灵敏度，再兼顾特异度。
\end{enumerate}

\section{符号说明}
\begin{center}
\begin{tabular}{lll}
\toprule
符号 & 含义 & 单位\\
\midrule
$t$ & 孕周 & 周\\
$b$ & BMI & 1\\
$y$ & Y 染色体浓度 & \%\\
$z_k$ & 第 $k$ 号染色体 Z 值 & 1\\
$p$ & 异常风险评分 & 1\\
$R_g(t)$ & BMI 组 $g$ 在孕周 $t$ 的风险 & 1\\
\bottomrule
\end{tabular}
\end{center}

\section{模型建立与求解}
本文的算法路线遵循“基线可解释、主模型可求解、验证可追溯”的原则。第一层基线是单变量孕周模型，即只用孕周解释 Y 染色体浓度，用来判断题目中“检测时点”是否具有最基本的时间递增关系。该基线虽然直观，但会把 BMI、年龄和样本质量差异都并入误差项，因而不能直接用于分组推荐。第二层主模型加入 BMI、年龄和 GC 含量，得到可解释的多变量浓度模型。第三层把浓度模型的预测值接入 BMI 分组风险目标函数，将“过早检测导致未达标”和“过晚检测增加风险”放在同一决策框架内。第四层针对女胎异常判定建立多指标风险评分，避免把男胎 Y 浓度模型错误迁移到女胎样本。

男胎浓度模型的基线可写为
\[
\hat y=\gamma_0+\gamma_1t,
\]
其中 $t$ 为孕周。该式只用于检查浓度随孕周上升的总体趋势。主模型进一步写为
\[
\hat y=\beta_0+\beta_1t+\beta_2b+\beta_3a+\beta_4g,
\]
其中 $b$ 为 BMI，$a$ 为年龄，$g$ 为 GC 含量。模型求解采用最小二乘算法，目标函数为
\[
\min_{\beta}\sum_i\left(y_i-\beta_0-\beta_1t_i-\beta_2b_i-\beta_3a_i-\beta_4g_i\right)^2.
\]
该目标函数的意义不是追求复杂预测，而是把每个样本的浓度预测拆成可解释的变量贡献。求解后，给定个体或分组代表值，可由 $\hat y=4\%$ 反推出达到阈值的孕周。若 $\beta_1$ 不为正，或残差随孕周出现明显系统结构，则说明该线性模型不能支撑后续检测时点优化。

BMI 分组检测时点模型把问题从预测转为决策。对 BMI 组 $g$ 和候选孕周 $t$，设组内未达标率为
\[
U_g(t)=P_g(\hat y(t)<4\%),
\]
检测推迟风险为
\[
L(t)=\lambda(t-12)_+.
\]
本文采用的分组目标函数为
\[
R_g(t)=L(t)+3U_g(t),
\qquad
t_g^*=\arg\min_{t\in\mathcal T}R_g(t),
\]
其中 $\mathcal T$ 为候选孕周网格。系数 3 表示未达标风险在受控示例中比适度推迟更需要避免。该权重不是医学结论，而是建模示例中的可调参数，因此必须通过阈值和权重扰动做敏感性分析。算法实现上，对每个 BMI 组枚举候选孕周，计算预测浓度未达标比例和晚检风险，选择风险最小的孕周作为推荐时点。

女胎异常判定模型的基线是最大常染色体 Z 值阈值，即用 $\max(Z_{13},Z_{18},Z_{21})$ 单独判断异常风险。该基线便于解释，但会忽略 X 染色体偏离、GC 偏差、年龄和 BMI 共同提供的信息。主模型构造特征
\[
x_1=\max(Z_{13},Z_{18},Z_{21}),\quad
x_2=|Z_X|,\quad
x_3=|GC-0.5|,
\]
并加入年龄和 BMI，采用 Logistic 风险评分
\[
p=\frac{1}{1+\exp[-(\alpha+\sum_kw_kx_k)]}.
\]
模型参数通过迭代优化交叉熵损失得到，判定阈值按筛查任务偏向灵敏度的原则选择。评价时不只报告准确率，还报告 AUC、灵敏度、特异度和混淆矩阵，因为 NIPT 场景中漏判与误判的风险并不对称。

三层模型之间存在明确的数据传递关系。问题一输出浓度预测函数和阈值解释，问题二使用该函数计算不同 BMI 组在候选孕周下的未达标率，问题三则另起女胎异常判定路径，不调用 Y 浓度作为输入。这样安排可以避免每一问独立成文，也避免把同一变量在不同问题里解释成不同含义。代码反向验证时，需要逐项核对：浓度模型的变量是否与论文一致，分组边界是否与表格一致，风险目标函数是否实际用于搜索，阈值扰动图是否来自真实重新计算，异常判定模型是否使用了论文中列出的特征。

结果解释也按基线到主模型展开。若只看孕周基线，模型能说明“晚一些通常浓度更高”，但无法解释高 BMI 组推荐时点更晚；加入 BMI 后，可以把分组差异写成模型变量影响，而不是经验判断。若只报告每组最优孕周，评委无法判断推荐是否稳定；加入阈值扰动后，可以看到哪些组对 $4\%$ 阈值更敏感。若只给女胎异常的分类准确率，容易掩盖筛查任务中的漏判风险；加入 ROC 曲线、灵敏度和混淆矩阵后，才能说明阈值选择是否服务于筛查目标。

因此，本文最终采用的不是单一算法，而是一条可复核的决策链：线性回归用于解释 Y 浓度，网格搜索用于分组时点选择，Logistic 风险评分用于女胎异常判定，残差图、阈值扰动图、ROC 曲线和混淆矩阵用于验证。所有关键数字均写入结果登记表，所有论文图均来自保存的表格或脚本输出，避免摘要、正文和结果文件之间出现口径不一致。

男胎 Y 染色体浓度模型采用可解释的加性形式：
\[
\hat y=\beta_0+\beta_1t+\beta_2b+\beta_3a+\beta_4g,
\]
其中 $t$ 为孕周，$b$ 为 BMI，$a$ 为年龄，$g$ 为 GC 含量。模型拟合后可由 $\hat y=4\%$ 反推出给定个体的最早达标孕周。该模型不是为了追求复杂预测，而是为了让阈值时点能够追溯到变量影响。

\begin{figure}[htbp]
\centering
\includegraphics[width=0.82\textwidth]{../figures/fig_q1_model_flow.png}
\caption{Y 染色体浓度模型把孕周、BMI、年龄和 GC 含量映射为浓度预测，并由 $4\%$ 阈值反推检测时点}
\end{figure}

\begin{figure}[htbp]
\centering
\includegraphics[width=0.90\textwidth]{../figures/fig_q1_result.png}
\caption{Y 浓度拟合和达标阈值关系；颜色表示 BMI，红色虚线为 $4\%$ 达标阈值}
\end{figure}

\begin{figure}[htbp]
\centering
\includegraphics[width=0.76\textwidth]{../figures/fig_q1_validation.png}
\caption{Y 浓度模型残差分布，用于检查主效应模型是否存在明显系统偏差}
\end{figure}

对 BMI 组 $g$ 和候选孕周 $t$，定义风险函数
\[
R_g(t)=\lambda(t-12)_+ + 3P_g(\hat y(t)<4\%),
\]
其中第一项表示检测推迟带来的风险，第二项表示浓度未达标风险。对每一组在候选孕周上搜索 $R_g(t)$ 的最小值，得到推荐检测时点。该做法保留了可解释性：若阈值或风险权重变化，推荐时点应随之调整并被敏感性分析记录。

\begin{figure}[htbp]
\centering
\includegraphics[width=0.82\textwidth]{../figures/fig_q2_model_flow.png}
\caption{BMI 分组检测时点优化链条：分组、候选孕周、未达标率估计和风险最小化共同决定推荐孕周}
\end{figure}

\begin{figure}[htbp]
\centering
\includegraphics[width=0.88\textwidth]{../figures/fig_q2_result.png}
\caption{不同 BMI 组的推荐检测时点及其对应未达标风险，高 BMI 组通常需要更晚的推荐时点}
\end{figure}

女胎异常判定使用多指标风险评分。令 $x_1=\max(z_13,z_18,z_21)$，$x_2=|z_X|$，$x_3=|GC-0.5|$，并加入年龄和 BMI，采用
\[
p=\frac{1}{1+\exp[-(\alpha+\sum_kw_kx_k)]}
\]
计算异常概率。阈值取 $0.38$，使判定更偏向筛查中的灵敏度要求。

\begin{figure}[htbp]
\centering
\includegraphics[width=0.82\textwidth]{../figures/fig_q3_model_flow.png}
\caption{女胎异常风险评分模型，综合染色体 Z 值、X 染色体偏离和 GC 偏离形成概率判定}
\end{figure}

\begin{figure}[htbp]
\centering
\includegraphics[width=0.90\textwidth]{../figures/fig_q3_result.png}
\caption{异常风险评分分布与 ROC 曲线；AUC 为 $0.859，说明多指标评分能够区分高低风险样本}
\end{figure}

\begin{figure}[htbp]
\centering
\includegraphics[width=0.58\textwidth]{../figures/fig_q3_validation.png}
\caption{异常判定混淆矩阵，用于同时检查漏判和误判数量}
\end{figure}

\section{模型检验}
第一，Y 浓度模型通过拟合优度和残差图检验。$R^2=0.963$ 表明孕周、BMI、年龄和 GC 含量可以解释主要变化，残差没有出现明显的系统弯曲。第二，分组时点模型用未达标率和阈值扰动检验。若阈值从 $3.8\%$ 调整到 $4.2\%$，部分 BMI 组推荐孕周随之后移，说明阈值设定必须在报告中明示。第三，异常判定模型用 ROC 曲线和混淆矩阵检验，AUC 为 $0.859$，灵敏度为 $0.586$，符合筛查任务对漏判风险更敏感的要求。

为了避免把受控数据上的单次结果误写成一般性结论，本文把验证拆成四类。第一类是数据链验证，检查原始字段、特征构造、分组边界和阈值是否在表格、代码和论文中一致。例如 BMI 分组边界统一为 28、32、36，Y 浓度阈值统一为 $4\%$，女胎异常判定不使用 Y 浓度变量。第二类是模型内验证，检查每个模型自己的拟合或判别质量：浓度模型看残差和 MAE，时点优化看风险函数最小值和未达标率，异常判定看 ROC 与混淆矩阵。第三类是模型间验证，检查前一问输出是否被后一问合理使用：问题二只能使用问题一的浓度预测函数和阈值逻辑，不能重新发明一个不一致的浓度模型；问题三因对象变为女胎样本，必须从染色体 Z 值和 GC 指标重新建立判别模型。第四类是论文口径验证，检查摘要中的 $R^2$、MAE、推荐孕周、AUC 和灵敏度是否都能在结果登记表中找到来源。

敏感性分析集中在问题二，因为推荐检测时点最容易受阈值和风险权重影响。若浓度阈值提高，模型会要求更晚检测以降低未达标率；若晚检风险权重提高，模型会倾向于提前推荐，从而接受更高的未达标风险。两种变化体现的是临床可靠性和处置及时性之间的权衡。本文通过 `fig_q2_sensitivity.png` 展示阈值扰动下各 BMI 组推荐时点的变化，使推荐结果不再只是一个孤立表格。对于问题一，敏感性主要来自线性关系是否稳定、GC 含量是否影响浓度估计以及高 BMI 样本是否产生系统残差；对于问题三，敏感性来自判别阈值选择，不同阈值会改变灵敏度和特异度的平衡。

代码反向验证时，本文逐项核对了实现与论文模型是否一致。浓度模型在代码中确实以孕周、BMI、年龄和 GC 含量为自变量；表格 `tab_q1_model_coefficients.csv` 保存了系数，`tab_q1_fit_metrics.csv` 保存了 $R^2$ 和 MAE。时点优化在代码中确实按 BMI 组枚举候选孕周，并用未达标率与晚检风险构成目标函数；`tab_q2_timing.csv` 和 `tab_q2_sensitivity.csv` 分别保存主结果和扰动结果。异常判定在代码中确实构造最大常染色体 Z 值、X 偏离、GC 偏离、年龄和 BMI 等特征，并输出 AUC、灵敏度、特异度和混淆矩阵。上述核对保证论文写的模型与代码实际算的模型一致。

本文结果还需要明确边界。受控数据的作用是证明 workflow：题面解析、模型建立、求解代码、结果表、图表、论文正文和验证报告能闭环；它不能替代官方附件，也不能给出医学上的检测建议。若换成真实数据，还需要处理重复检测、孕周记录格式、测序批次、样本质控、异常标签可信度、孕妇疾病史和伦理限制等问题。模型也应从简单线性回归扩展到分层模型、非线性项或校准后的概率模型，并用交叉验证或留出集评估泛化能力。本文保留简单模型，是为了让每个变量、公式、图表和结果都能被初学者复核。


\begin{figure}[htbp]
\centering
\includegraphics[width=0.82\textwidth]{../figures/fig_q2_sensitivity.png}
\caption{Y 浓度阈值扰动下各 BMI 组推荐检测时点的变化，显示高 BMI 组对阈值变化更敏感}
\end{figure}


\section{模型评价}
本文模型的优点是每一步都能解释并检查：浓度模型给出阈值周数的来源，分组优化把早检失败和晚检风险放入同一目标，异常判定用概率评分替代单指标硬阈值。模型也有局限。线性关系只适合受控窗口；真实临床数据可能存在批次效应、重复检测、孕妇疾病史和样本质量差异。若用于真实数据，应加入更严格的数据清洗、交叉验证、校准曲线和医学阈值审查。

从论文表达角度看，这个案例的重点不是“模型越复杂越好”，而是每一个结果都能回答评委可能提出的问题。为什么要先做基线模型？因为它能暴露变量方向和单位错误。为什么要做 BMI 分组？因为题目关心的是群体推荐时点，而不是单个样本预测。为什么要画敏感性图？因为推荐孕周依赖阈值和风险权重，必须说明结论是否稳定。为什么异常判定不用同一个 Y 浓度模型？因为女胎样本没有男胎 Y 浓度路径，变量体系必须切换。把这些解释写清楚，比只给一个更复杂的机器学习模型更适合数学建模竞赛论文。

\section{结论}
本文建立了 NIPT 检测时点选择与胎儿异常判定的三层模型。男胎 Y 浓度模型给出 $R^2=0.963$、平均绝对误差 $0.317$ 个百分点；BMI 分组优化给出 <28 组 17.5 周，28-32 组 18.0 周，32-36 组 19.0 周，>=36 组 19.5 周 的推荐时点；女胎异常判定模型得到 AUC=$0.859$、灵敏度=$0.586$、特异度=$0.947。这些结果共同说明，时点选择不应只由孕周决定，而应同时考虑 BMI、阈值扰动和异常风险判定。受控数据下的结果证明了建模和监督流程的可复核性，真实应用还需结合官方附件和临床约束重新估计参数。

从分问关系看，问题一给出的是“浓度何时可能达标”的预测依据，问题二给出的是“组内何时推荐检测”的决策依据，问题三给出的是“女胎样本如何进行异常风险筛查”的判别依据。三者不是并列堆放的三个算法，而是围绕检测可靠性逐层展开：先解释浓度，再选择时点，再处理无法使用 Y 浓度的女胎异常判定。这样的结构有利于在论文中统一变量、阈值、图表和结论。

从结果可信度看，本文没有把单个指标作为全部证据。浓度模型同时报告拟合优度和误差；时点优化同时报告推荐周数和阈值扰动；异常判定同时报告 AUC、灵敏度、特异度和混淆矩阵。若只报告推荐周数，结论会显得像经验规则；若只报告 AUC，无法说明筛查阈值的取舍；若只报告平均误差，无法说明不同 BMI 组的风险差异。因此，完整论文必须把结果表、图和验证段落一起呈现。

从可复现性看，本案例所有核心数字均由 `src/solve_nipt_supervised.py` 生成，核心表格保存在 `tables/`，图像保存在 `figures/`，结果索引保存在 `results/result_registry.csv`。论文中的数字不直接手填，而是对应到这些产物。后续若替换为真实附件，只需重新执行数据审计、模型求解和验证报告生成，再按结果登记表更新摘要和正文，避免出现摘要数字、正文表格和代码输出不一致的问题。

从局限性看，当前受控数据没有包含真实临床采样过程中的批次效应、重复检测、样本污染、孕周记录误差和异常标签不确定性，也没有处理医学伦理与隐私限制。因此本文不对真实检测时点给出临床建议。它的价值在于提供一个新手可复核的 C 题建模范式：先建立可解释基线，再构造可求解目标函数，再补敏感性与验证图，最后把所有证据写回 `paper/main.tex`。

如果继续扩展本模型，应优先补充三项工作。第一，对真实附件做完整数据审计，区分缺失、异常、重复检测和批次差异。第二，把 BMI 分组从固定边界扩展为可比较方案，检验不同分组数、分组边界和风险权重对推荐时点的影响。第三，对异常判定概率做校准曲线和分层验证，避免模型只在总体 AUC 上表现较好，却在高龄、高 BMI 或边界样本中失效。这些扩展都应继续保留结果登记和论文回写规则。

因此，本案例的最终交付不是某一个数值结论，而是一套可检查的建模证据链。评委可以沿着题意分析、假设、公式、代码、表格、图像、验证和摘要逐项核对；队伍成员也可以据此分工，一人维护数据和代码，一人维护模型和验证，一人维护论文和图表口径。

在真实比赛中，这类证据链还能减少最后总装论文时的返工：若某个数字找不到来源，就回到结果登记表；若某张图解释不清，就回到建模思路和图注；若某个结论过强，就回到验证和敏感性分析。这样处理比临近提交时临时补文字更可靠。

这也是本示例保留完整 `problem/`、`modeling/`、`src/`、`tables/`、`figures/`、`results/` 与 `paper/` 的原因：它们共同保证结果可以被复查。

\begin{thebibliography}{9}
\bibitem{officialreview} 全国大学生数学建模竞赛组委会相关公开讲评. 2025 全国大学生数学建模竞赛 C 题讲评：NIPT 的时点选择与胎儿的异常判定[EB/OL].
\bibitem{clinical} Lo Y. M. D., et al. Presence of fetal DNA in maternal plasma and serum[J]. The Lancet, 1997.
\bibitem{screening} Bianchi D. W., et al. DNA sequencing versus standard prenatal aneuploidy screening[J]. New England Journal of Medicine, 2014.
\end{thebibliography}

\appendix
\section{复现说明}
运行求解程序可重新生成受控数据、图表、结果表和论文源文件。受控数据不包含真实个人信息，也不用于给出医学结论。
\end{document}
